% =====================================================================
%  M2ISDA 2027 — Modele de resume / Abstract template
%  2e Conference internationale de modelisation mathematique
%  et d'innovations en sciences des donnees avancees
%  19 - 21 janvier 2027, Universite de Djibouti
%  Envoi du resume : lamo@univ.edu.dj   —   avant le 5 janvier 2027
%  Longueur : une page maximum, references comprises.
% =====================================================================
\documentclass[11pt,a4paper]{article}

\usepackage[utf8]{inputenc}
\usepackage[T1]{fontenc}
\usepackage[french,english]{babel}   % langue principale en dernier
\usepackage{amsmath,amssymb}
\usepackage[margin=2.5cm]{geometry}
\usepackage{graphicx}
\pagestyle{empty}

\begin{document}

% ---------------------------------------------------------------- TITRE
\begin{center}
  {\Large\bfseries Titre de la communication\par}
  \vspace{4mm}

  % Auteur presentant : souligner le nom
  \underline{Prenom Nom}$^{1}$, Prenom Nom$^{2}$, Prenom Nom$^{1,3}$
  \vspace{2mm}

  {\small
   $^{1}$ Laboratoire, Universite, Ville, Pays\\
   $^{2}$ Laboratoire, Universite, Ville, Pays\\
   $^{3}$ Institution, Ville, Pays\\
   \texttt{auteur.correspondant@exemple.org}
  }
\end{center}
\vspace{2mm}
\hrule
\vspace{4mm}

% ------------------------------------------------------------- RESUME
\noindent\textbf{Resume.}
Exposez en un paragraphe le probleme traite, le modele ou la methode
proposee, les donnees utilisees et les principaux resultats obtenus.
Les equations sont admises, par exemple
\begin{equation*}
  \frac{\partial u}{\partial t} = \Delta u + f(u), \qquad u(0)=u_{0},
\end{equation*}
ainsi qu'une figure ou un tableau si la place le permet.
Terminez par la portee des resultats pour la discipline ou pour
la decision publique.

\vspace{3mm}
\noindent\textbf{Mots-cles :} mot-cle 1, mot-cle 2, mot-cle 3, mot-cle 4.

\vspace{3mm}
\noindent\textbf{Axe thematique :} indiquer l'un des huit axes de la conference.

\vspace{1mm}
\noindent\textbf{Type de communication souhaite :} orale $\square$ \quad poster $\square$

% --------------------------------------------------------- REFERENCES
\vspace{4mm}
\begin{thebibliography}{9}\footnotesize
\bibitem{ref1} Nom, P., Nom, P. (2026).
  Titre de l'article. \emph{Nom de la revue}, 12(3), 145--167.
\bibitem{ref2} Nom, P. (2024).
  \emph{Titre de l'ouvrage}. Editeur, Ville.
\end{thebibliography}

\end{document}
